\documentclass{article}
\usepackage[T1]{fontenc}
\usepackage{lmodern}
\usepackage{graphicx}
\usepackage{amsmath,amssymb}
\usepackage[a4paper,margin=2cm]{geometry}
\usepackage[absolute,overlay]{textpos}
\setlength{\TPHorizModule}{1mm}
\setlength{\TPVertModule}{1mm}
\usepackage[indonesian]{babel}
\usepackage{hyperref}
\setlength{\emergencystretch}{2em}
% unit: Halaman 80

\begin{document}

% === Halaman 80 (llm) ===

\begin{center}
\Large \textbf{The Salsa20 family of stream ciphers}
\vspace{1cm}
\large Daniel J. Bernstein *
\small Department of Mathematics, Statistics, and Computer Science (M/C 249)\
The University of Illinois at Chicago\
Chicago, IL 60607-7045\
\texttt{snuffle66@box.cr.yp.to}
\end{center}

\noindent \textbf{Abstract.} Salsa20 is a family of 256-bit stream ciphers designed in 2005 and submitted to eSTREAM, the ECRYPT Stream Cipher Project. Salsa20 has progressed to the third round of eSTREAM without any changes. The 20-round stream cipher Salsa20/20 is consistently faster than AES and is recommended by the designer for typical cryptographic applications. The reduced-round ciphers Salsa20/12 and Salsa20/8 are among the fastest 256-bit stream ciphers available and are recommended for applications where speed is more important than confidence. The fastest known attacks use $\approx 2^{153}$ simple operations against Salsa20/7, $\approx 2^{249}$ simple operations against Salsa20/8, and $\approx 2^{255}$ simple operations against Salsa20/9, Salsa20/10, etc. In this paper, the Salsa20 designer presents Salsa20 and discusses the decisions made in the Salsa20 design.

\section{Introduction}

A sender and receiver share a short secret key. They use the secret key to encrypt a series of messages. A message could be short, just a few bytes, but it could be

\end{document}
